\documentclass[11pt]{amsart}
\usepackage[T1]{fontenc}
\usepackage{lmodern,amsmath,amssymb,mathtools,microtype,booktabs}
\usepackage[margin=1.15in]{geometry}
\usepackage[hidelinks]{hyperref}
\newtheorem{theorem}{Theorem}[section]
\newtheorem{proposition}[theorem]{Proposition}
\newtheorem{lemma}[theorem]{Lemma}
\newtheorem{corollary}[theorem]{Corollary}
\theoremstyle{definition}\newtheorem{definition}[theorem]{Definition}
\theoremstyle{remark}\newtheorem{remark}[theorem]{Remark}
\newcommand{\R}{\mathbb R}
\newcommand{\Z}{\mathbb Z}
\newcommand{\dd}{\,\mathrm d}
\newcommand{\Prob}{\mathbb P}
\DeclareMathOperator{\covol}{covol}
\DeclareMathOperator{\dist}{dist}
\DeclareMathOperator{\area}{area}
\numberwithin{equation}{section}
\title[Intrinsic marked boundary laws]{Intrinsic marked boundary laws and rigidity of periodic dispersing billiards}
\author{Qian Qi}
\date{September 28, 2026}
\subjclass[2020]{37D50, 53A04, 70H09, 62G05}
\keywords{Dispersing billiards, intrinsic endpoint laws, local scattering,
period-lattice reconstruction, unlabelled channel data, conditional stability}
\hypersetup{pdftitle={Intrinsic marked boundary laws and rigidity of periodic dispersing billiards},pdfauthor={Qian Qi}}
\begin{document}
\begin{abstract}
Two strictly active endpoint-arclength density functions of an isolated
return channel determine its stationary two-flight action, both smooth
reflecting arcs, their gap, and the free area. We formulate the exact
relation with local travel-time and scattering data and give a conditional
Lipschitz inverse. For analytic periodic tables with asymmetric,
pairwise noncongruent obstacle representatives, an unlabelled collection
of these density pairs determines the incidence graph and its geometric
cycle displacements: no obstacle identifications, cross-channel contact
registrations, or integer lattice-copy labels are supplied. The recovered
free area and obstacle areas determine the period covolume. Two independent
cycle displacements then give a finite, explicit list of possible lattices;
when their determinant equals the recovered covolume, the entire table
is unique up to isometry. We also construct physical finite-index
ambiguities. On compact classes with uniform analytic strips, geometric
and symmetry margins, shape separation, and onset brackets, finite growing
endpoint histograms give a fixed-exponent conditional reconstruction,
including stable recovery of the discrete incidence and primitive period
structure. The observations remain marked by local contact origins and
selected itineraries. Earlier moment, count-fiber, registered-network,
and relative-law results are retained in the accompanying supplementary
volumes. No finite-horizon assumption is imposed.
\end{abstract}
\maketitle
\section{The observation and the main results}\label{sec:setting}
Near the onset of a prescribed billiard itinerary, phase volume contains
a residual-time factor as well as a collision-section flux. Keeping
intrinsic endpoint positions rather than integrating them out makes
this structure visible. Two density functions recover the local action
and its volume normalization. The normalization has a second use:
after analytic boundary recovery it determines the covolume of the
unknown period lattice. This connects the local observation to a
finite-index global reconstruction without supplied lattice coordinates.

\subsection{Local records}
A periodic dispersing table is the complement of
\[
             \bigcup_{i=1}^{r}\ \bigcup_{\lambda\in\Lambda}
                         (C_i+\lambda)\subset\R^2,
\]
where $\Lambda$ is a full-rank lattice and the $C_i$ are compact strictly
convex bodies with smooth positively curved boundaries. All translates
are disjoint with positive separation. There is one representative per
obstacle orbit, and the free area is
\begin{equation}\label{eq:volume}
             A=\covol\Lambda-\sum_{i=1}^r\area(C_i)>0.
\end{equation}
Initial conditions have speed one and normalized phase measure
$\dd q\dd\varphi/(2\pi A)$. A selected channel is an isolated,
unobstructed common normal between two obstacle copies. Small facing
patches and a sufficiently short onset collar select exactly the three
successive impacts source--opposite--source.

On success record the first and last signed source arclengths $s,t$
from its normal contact. The corresponding sign on the opposite arc
increases in the same transverse direction at the normal orbit.
Only the source arclengths are observed. Failure is a separate outcome.
At absolute time $T$ the subprobability measure is
\[
 \mu_T(B)=\Prob\{\text{the selected three-impact event by }T,
                                      \ (s,t)\in B\}.
\]
It is not conditioned on success. A local record consists of the two
absolute times $T_2>T_1$ and their measures, on a common strictly active
rectangle, modulo one simultaneous reversal $(s,t)\mapsto(-s,-t)$.
The sign choice is the same at both times. Each record retains its own
contact origin, itinerary selection and grouping of its observations.
No Cartesian position, direction, gap or area is supplied.

The phrase \emph{two density functions} will always mean function-valued
data. Their finite counterparts below are two increasingly fine
histograms, with a number of cells tending to infinity. Neither is a
pair of scalar observations. The local pilot only chooses active times;
it need not estimate the onset to vanishing error.

\begin{theorem}[The smooth local inverse]\label{thm:local-main}
For one selected channel, two strictly active endpoint density functions
$f_1,f_2$ recover its action and area by
\begin{equation}\label{eq:ratio-main}
 W=\frac{T_1f_2-T_2f_1}{f_2-f_1},\qquad
 A=\frac{(T_2-T_1)(-W_{st})}{2\pi(f_2-f_1)}.
\end{equation}
The diagonal gradient and Hessian of $W$ determine both smooth boundary
arcs in their relative placement, up to a common Euclidean isometry.
The gap is $W(0,0)/2$. Under the explicit physical $C^3$, active,
twist and time margins in Proposition~\ref{prop:local-stable}, the
inverse is Lipschitz from the two $C^2$ densities to the intrinsic
$C^2$ arcs, gap and area. Analyticity is unnecessary for this local result.
\end{theorem}

\subsection{An unlabelled collection and a covolume test}
Call a body asymmetric if its Euclidean isometry group is trivial.
Call a table shape-separated if its distinct representatives $C_i$
are pairwise noncongruent, allowing reflections in this comparison.
Repeated observations of translates of the same representative are
allowed. These are conditions on the unknown table, not supplied
boundary shapes or cross-experiment correspondences.

An \emph{unlabelled collection} is a finite multiset of local records.
The pairing of the two windows within a record and the local return
orientation are retained. There is no supplied incidence graph, obstacle
label, cross-channel contact identification, relative sign, period basis,
or integer copy label. Equality allows a permutation of records and an
independent coherent reversal of each record. The underlying channels
are required to visit every obstacle orbit and to form a connected
incidence graph. Coverage and connectedness are properties of the
experiment; they are not inferred for an obstacle never visited.

For analytic boundaries each reconstructed arc determines the entire
boundary image. Congruence classes of the images recover the vertices
of the incidence graph. Matching along a spanning tree places one
representative of each class. A remaining edge gives a translation
$d_e$ between its placed target image and the chosen target representative.
These operations and their independence of all local sign choices are
proved in Section~\ref{sec:descent}. Put
\begin{equation}\label{eq:calibrated-volume}
 V=A+\sum_i\area(C_i),\qquad
 D=[d_{e_1}\ d_{e_2}],\qquad n=\frac{|\det D|}{V},
\end{equation}
where $d_{e_1},d_{e_2}$ are any two independent non-tree displacements.
They need not have been marked by integers.

\begin{theorem}[Unlabelled finite-index descent]\label{thm:descent-main}
Consider an analytic asymmetric, shape-separated table and a connected,
covering unlabelled collection with two independent reconstructed cycle
displacements. The collection determines the obstacle images, incidence
graph, tree placement, free area and period covolume. The number $n$
in \eqref{eq:calibrated-volume} is a positive integer. Every compatible
period lattice is in the following explicit finite list:
\begin{equation}\label{eq:hnf-main}
 \Lambda_B=DB^{-1}\Z^2,\qquad
 B=\begin{pmatrix}a&b\\0&d\end{pmatrix},\quad
 ad=n,\quad a,d>0,\quad 0\le b<d.
\end{equation}
Thus there are at most $\sigma_1(n)=\sum_{d\mid n}d$ compatible tables,
up to a common isometry and permutation of obstacle representatives.
Other cycle, separation and channel constraints can only remove candidates.
If $|\det D|=V$, the whole table and its unmarked lattice are unique.
All assertions use the two endpoint density functions per record, not
counts. No global integer marking or incidence identification is an input.
\end{theorem}

The determinant condition is a check on recovered geometric quantities,
not a supplied primitive pair of integer labels. It is stronger than
merely having two independent cycles. The finite-index qualification is
necessary: Proposition~\ref{prop:ambiguity} constructs, for every odd
prime $p$, $p-1$ nonisometric physical tables with identical selected
local records and calibrated index $p$. Theorem~\ref{thm:descent-main}
also has nonempty open primitive classes with any prescribed finite
number of distinct obstacle shapes.

\begin{theorem}[Finite observations on a primitive class]\label{thm:stat-main}
Fix a compact class of the preceding tables with uniform holomorphic
support-function strips, positive curvature, separation and channel
margins, bounded geometry, quantitative asymmetry and separation between
distinct shape classes. Each selected onset has a known bounded bracket;
the number $J$ of records is bounded. Suppose a reconstructed cycle pair
satisfies $|\det D|=V$. No identity of that pair is supplied to the estimator.

Let $0<\beta\le1$ and assume uniform $C^{2,\beta}$ active-density bounds.
A finite pilot followed by two growing endpoint histograms per record
gives, with probability at least $1-\delta$, whole-table error at most
\begin{equation}\label{eq:rate-main}
 C\left(\frac{\log(CJN/\delta)}{N}\right)^{\beta\theta/(2\beta+6)},
                       \qquad 0<\theta<1,
\end{equation}
using at most $CJN$ independent phase preparations, including failures.
The incidence classes and the primitive period structure are stable for
sufficiently large $N$. The exponent $\theta$, constants and the required
sample size depend on the entire prior class. The loss and all margins
are specified in Section~\ref{sec:stability}. This is a conditional upper
bound, not a minimax exponent or a claim about a count-only sensor.
\end{theorem}

The proof passes from histograms to actual arcs, then to analytic images,
and finally to the finite incidence and period data. It does not estimate
an increasing number of contact jets. The new global step uses the area
normalization in \eqref{eq:ratio-main}, which would be lost by conditioning
on successful trajectories. Symmetric tables, repeated congruent obstacle
representatives and nonprimitive collections retain their earlier marked
or finite-candidate results; they are not silently included in the new
unique unlabelled theorem.

\section{The local action and two reflecting arcs}\label{sec:local}
\subsection{The physical return branch}
Place a normal contact pair at $(0,0)$ and $(g,0)$ for the proof only.
In transverse graph coordinates its boundaries are
$(-\psi_0(y),y)$ and $(g+\psi_1(y),y)$, where
$\psi_b(0)=\psi_b'(0)=0$ and $\psi_b''(0)=\kappa_b>0$.
The broken length has a nondegenerate minimum in its intermediate
coordinate: its second derivative there is $2(g^{-1}+\kappa_o)>0$.
The implicit-function theorem gives a unique stationary intermediate
impact for endpoints in a smaller fixed box. Strict channel clearance
identifies this branch with the physical specular trajectory.

Use intrinsic unit-speed boundary parametrizations $x(s)$ and $y(u)$,
with both tangents in the positive transverse direction at zero, and put
\begin{equation}\label{eq:action}
 D(s,u)=|x(s)-y(u)|,\qquad
 W(s,t)=D(s,u(s,t))+D(t,u(s,t)).
\end{equation}
The stationary equation is $D_u(s,u)+D_u(t,u)=0$. On a sufficiently
small rectangle $W$ is smooth, symmetric in its two arguments, and has
positive dispersing twist $-W_{st}>0$. At the normal orbit,
\[
 W(0,0)=2g,\qquad W_{st}(0,0)=-\frac1{2g(1+g\kappa_o)}.
\]
These statements hold uniformly on compact physical families with the
stated geometry and the required fixed derivative bounds.

\begin{lemma}[The residual-time density]\label{lem:density}
On a sufficiently short onset collar the endpoint density is
\begin{equation}\label{eq:density}
 f_T(s,t)=\frac{-W_{st}(s,t)}{2\pi A}(T-W(s,t))_+.
\end{equation}
On a rectangle where $T_1-W>0$, the two density functions at
$T_2>T_1$ recover $W$ and $A$ by \eqref{eq:ratio-main}.
\end{lemma}
\begin{proof}
Let $p$ be the momentum conjugate to source arclength and $r$ the time
before its first impact. Collision-section phase volume is
$\dd s\dd p\dd r$. The generating identity $p=-W_s$ changes this to
$(-W_{st})\dd s\dd t\dd r$. The allowed residual interval is
$0<r<T-W(s,t)$. Choose the collar shorter than a common positive lower
bound for adjacent free flights. Then this interval cannot include a
preceding collision, and the time after the third impact cannot include
a fourth. The strict local minimum and patch isolation keep the active
sublevel in this physical chart. Conversely, every selected near-onset
trajectory has these coordinates. Integrating $r$ and normalizing by
$2\pi A$ proves \eqref{eq:density}.

On a strictly active rectangle set $w=-W_{st}/(2\pi A)>0$.
Then $f_j=w(T_j-W)$ and $f_2-f_1=(T_2-T_1)w$.
Eliminating $w$ proves the action formula, and its recovered mixed
second derivative proves the area formula. Equality of measures
uniquely determines their continuous density representatives.
\end{proof}

\begin{proposition}[Exact local information equivalence]\label{prop:lens}
On the physical class and a fixed strictly active rectangle, the two
endpoint density functions at known absolute times are equivalent to
$(W,A)$. In particular they determine the local canonical scattering
relation
\begin{equation}\label{eq:lens}
       (s,-W_s(s,t))\longmapsto(t,W_t(s,t)).
\end{equation}
Conversely, that relation, one absolute action value, and $A$ determine
the densities on the rectangle. No extension to a global exterior
travelling-time relation is asserted.
\end{proposition}
\begin{proof}
The forward and inverse formulas are Lemma~\ref{lem:density}.
The positive twist makes $t\mapsto-W_s(s,t)$ locally invertible.
Thus \eqref{eq:lens} is a canonical graph; its generating one-form is
$-p_-\dd s+p_+\dd t=\dd W$. On the connected rectangle it determines
$W$ up to a constant, fixed by one value. Substitution in
\eqref{eq:density} gives the converse. Here the given relation is assumed
physical, so its generating one-form is exact. Arbitrary positive
functions satisfying the quotient formula are not being declared
physical endpoint laws.
\end{proof}

This makes explicit both the richness and the origin of the data.
First derivatives of travel time recover endpoint directions in
scattering geometry; compare Noakes--Stoyanov~\cite[Lemma 4]{NS} and
Gurfinkel--Noakes--Stoyanov~\cite[Lemma 7]{GNS}. The affine elimination
is not a new abstract lens-rigidity principle. What remains to prove
here is an explicit inverse for two unknown reflecting arcs in their
intrinsic coordinates and its periodic descent with volume calibration.

\subsection{The curvature identities}
On the diagonal write $u(s)=u(s,s)$ and define
\begin{equation}\label{eq:diagonal}
 \ell(s)=\tfrac12W(s,s),\qquad a(s)=W_s(s,s),\qquad
 v(s)=\sqrt{1-a(s)^2}.
\end{equation}
Shrink the interval so that $v>0$. Let $T=x'$ and $N$ be the source
unit tangent and inward normal, with $T'=kN$ and $N'=-kT$.
The vector $e=(x-y)/\ell$ then satisfies $e=aT+vN$.

\begin{lemma}[Diagonal Schur complement]\label{lem:curvature}
The source curvature and the opposite curvature at the nearest point are
\begin{align}
 k(s)&=\frac{W_{ss}(s,s)-W_{st}(s,s)-v(s)^2/\ell(s)}{v(s)},
                                                    \label{eq:curvature-source}\\
 k_o(u(s))&=\frac1{\ell(s)}
       \left(\frac{v(s)^2}{-2\ell(s)W_{st}(s,s)}-1\right),
                                                    \label{eq:curvature-other}\\
 u'(s)&=\frac{v(s)}{1+\ell(s)k_o(u(s))}>0.          \label{eq:foot-speed}
\end{align}
\end{lemma}
\begin{proof}
At $s=t$ the stationary equation is $D_u=0$, selecting the nearest
point on the opposite strictly convex patch. Orient its tangent $T_o$
continuously so that $T\cdot T_o=v>0$. Direct differentiation gives
\[
 D_{ss}=v^2/\ell+kv,\qquad D_{su}=-v/\ell,\qquad
 D_{uu}=1/\ell+k_o.
\]
For the last identity, $y''=-k_oe$ at the nearest point; this is the
external, dispersing sign. The stationary elimination of the intermediate
variable gives
\begin{equation}\label{eq:schur}
 W_{ss}=D_{ss}-\frac{D_{su}^2}{2D_{uu}},\qquad
 W_{st}=-\frac{D_{su}^2}{2D_{uu}}
        =-\frac{v^2}{2\ell(1+\ell k_o)}.
\end{equation}
Subtracting the first two expressions gives
\eqref{eq:curvature-source}; solving the last gives
\eqref{eq:curvature-other}. Differentiation of $D_u(s,u(s))=0$
yields \eqref{eq:foot-speed}.
\end{proof}

\begin{proof}[Proof of Theorem~\ref{thm:local-main}, exact assertion]
Recover $W$ and $A$ from Lemma~\ref{lem:density}. Choose
$x(0)=0$, $T(0)=(0,1)$ and $N(0)=(-1,0)$ and solve the Frenet
system with the recovered $k$. This determines the source arc. The
opposite arc in the same plane is
\begin{equation}\label{eq:foot}
 y(u(s))=x(s)-\ell(s)\{a(s)T(s)+v(s)N(s)\}.
\end{equation}
By \eqref{eq:foot-speed} it is a regular arc containing
$y(0)=(g,0)$, where $g=W(0,0)/2$. The recovered objects are functions
on intervals, not merely their Taylor series. Simultaneous arclength
reversal reflects the entire pair, not its members separately. The
remaining choices are one translation and orthogonal transformation.
\end{proof}

\begin{proposition}[Conditional local Lipschitz inverse]\label{prop:local-stable}
Fix nested endpoint rectangles $Q_0\Subset Q$, a diagonal subinterval,
and positive constants $M,g_*,A_*,v_*,w_*,\tau_*,m_*$. Consider smooth
physical pairs for which, in their normal contact frames, both intrinsic
boundary parametrizations have $C^3$ norm at most $M$ on fixed slightly
larger arcs, $g_*,A_*\le g,A\le M$, $|T_j|\le M$,
$T_2-T_1\ge\tau_*$, $T_1-W\ge m_*$ on $Q$,
$-W_{st}\ge w_*$ on $Q$, $v\ge v_*$ on the diagonal, and
$\|f_j\|_{C^2(Q)}\le M$. The patches, isolation and adjacent free-flight
margins are uniform, as is positive curvature. For two such pairs at
the same times put
\[
        e=\max_{j=1,2}\|f_j-\widetilde f_j\|_{C^2(Q)}.
\]
After a permitted common reflection, their gaps, areas, and intrinsic
$C^2$ boundary arcs on smaller fixed intervals differ by at most $Ce$.
The constant depends on all the stated bounds and margins, not on a
contact-jet order. This is an estimate between physical density pairs,
not an inverse defined on every pair of $C^2$ functions.
\end{proposition}
\begin{proof}
The density difference is bounded below by
$\tau_*w_* /(2\pi M)$. Products and reciprocals are Lipschitz on the
corresponding bounded $C^2$ sets. Formula~\eqref{eq:ratio-main} gives
$C^2$ action error $Ce$; its area expression gives scalar error $Ce$.
The diagonal formulas give $Ce$ source-curvature error and opposite
curvature error as functions of $s$, since all their denominators have
positive margins. The Frenet integral equations and Gronwall's inequality
give $Ce$ error for the source arc through its second derivative.

For the opposite arc, differentiating \eqref{eq:foot} twice would
unnecessarily require a third derivative of $W$. Instead integrate
\eqref{eq:foot-speed}. Its speed has positive upper and lower bounds,
and its error is $Ce$. Consequently $u(s)$ and its inverse have error
$Ce$ on smaller intervals. The physical $C^3$ prior bounds the derivative
of opposite curvature in intrinsic arclength. Thus the two opposite
curvature functions differ by $Ce$ at the same $u$. Their Frenet systems,
with the recovered gap and common initial tangent, give the required
$C^2$ estimate. This uses exactly the physical prior stated in the
proposition and no derivative of an arbitrary noisy curvature estimate.
\end{proof}

\section{Incidence and period descent without global labels}\label{sec:descent}
\subsection{Complete images and their incidence}
For a convex body $C$ let $p_C$ be its support function after translation
to its Steiner point. Congruence includes orthogonal transformations of
either determinant. A strictly convex real-analytic boundary has a
real-analytic support function on the circle: positive curvature makes
the normal-angle map analytically invertible. Two such bodies sharing
an open boundary arc in one frame have equal support functions on an
open normal-angle interval, hence everywhere by the identity theorem.
Thus Theorem~\ref{thm:local-main} determines each entire analytic pair
in its relative placement. No smooth-jet identity is used for this step.

\begin{lemma}[Intrinsic incidence recovery]\label{lem:incidence}
For an asymmetric, shape-separated table, a connected covering unlabelled
collection determines its incidence multigraph and places one representative
body of each vertex along any chosen spanning tree, up to a common isometry.
Every non-tree edge determines a unique translation $d_e$ between its
target image and the chosen representative. Each $d_e$ belongs to the
unknown period lattice. The construction uses no integer edge labels.
\end{lemma}
\begin{proof}
Reconstruct the two complete images from each local record. Partition
all image occurrences by Euclidean congruence. Shape separation means
that two occurrences are in the same class precisely when their physical
bodies are translates of the same obstacle representative. Coverage makes
the classes exactly the vertices; the two image slots of each record
give its incident vertices, including loops and repeated edges.

If $C$ and $C'$ are congruent asymmetric bodies, there is exactly one
isometry carrying $C$ to $C'$. Indeed two such maps differ by a symmetry
of $C$. Place one root image arbitrarily. Along an outward tree edge,
match its source image to the already placed source representative and
apply the same isometry to the opposite image. This places the child.
For backward traversal interchange the two images; no additional
observation is required. A change of the local reversal representative
is cancelled by composing its unique matching map with that reversal.

On a non-tree edge match the source in the same way. Let $\widehat C_w$
be the target image and $C_w$ the chosen target representative. In a
physical table they differ by a period translation. This translation is
unique since a nonempty compact set has no nonzero translational symmetry.
It is computed from Steiner points:
\begin{equation}\label{eq:cycle-displacement}
 d_e=s(\widehat C_w)-s(C_w),\qquad \widehat C_w=C_w+d_e.
\end{equation}
The second equality is a consistency condition, not an extra matching
choice. Tree placement chooses physical copies of the true obstacles;
therefore $d_e\in\Lambda$. All output changes under root placement by
one common orthogonal transformation and translation only.
\end{proof}

The graph is recovered from shapes, not from a comparison of curvature
values at isolated contacts. Pairwise noncongruence is essential for
this particular incidence argument. It is not an assumption that a
shape catalogue has been supplied to the observer. On the exact class,
congruence is an equality test on recovered analytic functions; no
polynomial-time test of such equality is asserted.

\subsection{The normalization determines the covolume}
Every record supplies the same $A$ in \eqref{eq:ratio-main}. Since the
incidence classes contain exactly one obstacle representative per period
cell, the recovered images give
\[
             V=A+\sum_i\area(C_i)=\covol\Lambda.
\]
This identity would not hold with a sum over image occurrences, which
would count repeated observations several times. It is also why the
shape-separated hypothesis cannot simply be discarded while removing
obstacle labels. Areas can be computed intrinsically from the images;
for example
\begin{equation}\label{eq:support-area}
 \area(C)=\frac12\int_0^{2\pi}(p_C^2-(p_C')^2)\dd\theta.
\end{equation}
The formula is unchanged by translating $C$.

\begin{lemma}[Finite-index lattice classification]\label{lem:overlattice}
Let $D$ be an invertible real $2\times2$ matrix and $V>0$. The lattices
$\Lambda$ containing $D\Z^2$ and having covolume $V$ form the list
\eqref{eq:hnf-main}, provided $n=|\det D|/V$ is a positive integer.
Otherwise there is no such lattice. The list has exactly $\sigma_1(n)$
distinct lattices. An additional vector $d_e$ belongs to $\Lambda_B$
if and only if $BD^{-1}d_e\in\Z^2$.
\end{lemma}
\begin{proof}
The subgroup $D\Z^2\subset\Lambda$ has index $n$, by the ratio of
fundamental parallelogram areas. Normalize by $D^{-1}$ and write
$L=D^{-1}\Lambda\supset\Z^2$. Its dual
$L^*=\{z:z\cdot L\subset\Z\}$ is a sublattice of $\Z^2$ of
index $n$. Conversely every such sublattice is the dual of a unique $L$.

For completeness, a sublattice $H\subset\Z^2$ of finite index has
a unique row basis $(a,b),(0,d)$ with $a,d>0$ and $0\le b<d$.
Its projection on the first coordinate is $a\Z$, its vertical intersection
is $\{0\}\times d\Z$, and subtraction of vertical multiples gives the
unique representative $b$. Every element then subtracts to zero using
these two rows. Its index is $ad$. Thus $H$ is the row span of the
matrix $B$ in \eqref{eq:hnf-main}, and its dual is $B^{-1}\Z^2$.
For each $d\mid n$ there are $d$ choices of $b$. This proves the count,
uniqueness of the list, and the stated membership test.
\end{proof}

\begin{proof}[Proof of Theorem~\ref{thm:descent-main}]
Lemma~\ref{lem:incidence} recovers all representative images and cycle
displacements. The normalization recovers $V$. Two independent $d_e$
generate a sublattice $D\Z^2$ of the true period lattice, so
Lemma~\ref{lem:overlattice} applies. For each listed lattice use the
already placed representative bodies. Test all remaining cycle vectors,
disjointness of translates, area and the selected physical channels.
A true table producing the records appears on the list and passes every
test. There is no additional continuous placement parameter. The list
therefore bounds the number of compatible tables. If $n=1$ it consists
only of $D\Z^2$.

The outcome is independent of the initial local reflections and root
placement by Lemma~\ref{lem:incidence}. A different spanning tree changes
the chosen copies and the cycle generators but not any physical table
passing the tests. One can enumerate the finitely many tree and cycle-pair
choices to find a determinant equal to $V$; the observer need not be told
which pair has this property. This proves the unique case as well.
\end{proof}

\begin{remark}[The unmarked period lattice]
No preferred basis of $\Lambda$ is recovered or required. In the primitive
case $D$ is one recovered basis, up to the usual unimodular change. On the
shape-separated class $\Lambda$ is in fact the full translational period
group: a translation preserving the table must preserve each shape class,
and hence sends $C_i$ to $C_i+\lambda$ for some $\lambda\in\Lambda$.
Compactness forces that translation to be $\lambda$. Thus the conclusion
does not determine an artificially chosen nonminimal covering torus.
\end{remark}

\subsection{Physical primitive classes and finite-index ambiguity}
\begin{proposition}[Nonempty open primitive classes]\label{prop:primitive}
For every positive integer $r$ there are nonempty open analytic families
with $r$ pairwise noncongruent asymmetric obstacle representatives and
a connected covering collection satisfying $|\det D|=V$.
\end{proposition}
\begin{proof}
Use small strictly convex bodies with centered support functions
\[
 p_i(\theta)=r_i\{1+\varepsilon\cos(2\theta)
                    +\varepsilon\cos(3\theta+\pi/4)\},
                \qquad 0<|\varepsilon|<1/11,
\]
where the positive $r_i$ are distinct. The curvature radius $p_i+p_i''$
is positive. A rotation preserving both Fourier modes has angle zero
modulo $2\pi$. A reflection about angle $a$ would require both
$2a\in\pi\Z$ and $3a+\pi/4\in\pi\Z$, which is impossible.
Distinct scales are noncongruent.

Take a rectangular lattice with periods $(L,0),(0,M)$. Place one small
body at the origin and the others at generic points strictly inside the
cell. Choose a star joining the root to these points, and include the
root's horizontal and vertical neighboring-copy channels. Generic points
ensure that no other center lies on any chosen segment, including relevant
nearby translated centers. Since there are only finitely many such
segments and nearby centers, sufficiently small bodies leave positive
clearance. The normal channels persist by the nondegenerate stationary
minimum. The star is a spanning tree, and the two root loops yield exactly
$(L,0)$ and $(0,M)$, whose determinant equals the period covolume.
Positive clearance, convexity, asymmetry and pairwise noncongruence persist
in a sufficiently small analytic neighborhood. The loop labels need not
be observed: their geometric displacements remain the two periods.
\end{proof}

\begin{proposition}[Physical finite-index ambiguity]\label{prop:ambiguity}
For each odd prime $p$ there are $p-1$ nonisometric analytic periodic
tables with identical unlabelled two-density records, identical free area,
a connected covering graph and two independent reconstructed displacements,
for which the calibrated index is $p$.
\end{proposition}
\begin{proof}
Fix one asymmetric analytic strictly convex body $C$ contained in a disk
of radius $r$, and choose $L,M>8pr$. Put $d_1=(L,0)$, $d_2=(0,M)$
and, for $j=1,\ldots,p-1$, set
\[
 \Lambda_j=\Z d_1+\Z d_2+
                      \Z\frac{d_1+j d_2}{p}.
\]
These are distinct superlattices of $\Z d_1+\Z d_2$ of index $p$;
each has covolume $LM/p$. Their nonzero vectors have length at least
$\min(L,M)/p$, so the copies of $C$ are disjoint. All tables have free
area $A=LM/p-\area(C)>0$.

Observe only the channels from $C$ to $C+d_1$ and from $C$ to $C+d_2$.
A center in $\Lambda_j$ on the horizontal axis is an integral multiple
of $d_1$, because $j$ is invertible modulo $p$. A center on the vertical
axis is an integral multiple of $d_2$. Off the horizontal axis its
vertical distance is at least $M/p$, and off the vertical axis its
horizontal distance is at least $L/p$. Thus the two selected pair
corridors, lying within radius $r$ of their axis segments, are separated
from every other body by a positive common margin. Their normal-channel
geometries are independent of $j$. On sufficiently small common patches
and time collars, Lemma~\ref{lem:density} gives exactly the same densities:
both $W$ and $A$ agree. This is equality of physical laws, not merely
of finite jets. The recovered graph has one vertex and two loops, and
$|\det[d_1\ d_2]|/V=p$.

Finally an isometry between two tables sends a connected obstacle component
to a translate of $C$. Its orthogonal part must preserve the centered body,
and is therefore the identity. A translation conjugates neither period
group to a different one. Since $\Lambda_j$ are distinct and are the full
period groups, the tables are nonisometric.
\end{proof}

The proposition locates the finite obstruction rather than replacing the
rigidity theorem by nonidentifiability. The primitive theorem gives unique
reconstruction on an open physical class; the general theorem exhausts
all remaining period choices by a finite arithmetic list.

\section{Conditional finite-data reconstruction}\label{sec:stability}
\subsection{The prior class and the loss}
Fix a compact class of physical presentations with at most $J$ records
and at most $2J$ obstacle representatives. All selected channels have
uniform positive separation, isolation and adjacent free-flight margins;
gaps, curvature, free area and period covolume have positive lower bounds.
Diameters, gaps, areas and some lattice bases and their inverses have
fixed upper bounds. Each selected onset belongs to a known bounded bracket,
and queries in a fixed collar beyond its upper endpoint are allowed.
These are local itinerary brackets, not exact onsets or hidden contacts.

Steiner-centered support functions extend holomorphically to
$|\operatorname{Im}z|<\rho$ with a common bound $M$. For constants
$\eta,\nu>0$ require
\begin{align}
 \|p_i(\cdot+\phi)-p_i\|_2
   &\ge\eta\dist(\phi,2\pi\Z),                 \label{eq:rot-margin}\\
 \|p_i(-\cdot+\phi)-p_i\|_2&\ge\eta,         \label{eq:ref-margin}\\
 \inf_{O\in O(2)}\|p_i-O\cdot p_j\|_2&\ge\nu\quad(i\ne j).
                                                        \label{eq:shape-margin}
\end{align}
The norms are on the real normal-angle circle. Every fixed finite
asymmetric, shape-separated family admits such margins. Near zero rotation,
the first follows from $\|p_i'\|_2>0$; elsewhere and for reflections or
distinct shapes, compactness of $O(2)$ and absence of equality give the
positive bounds. The bounds persist in a small compact neighborhood.
They are quantitative priors, not a supplied shape catalogue.

Uniform analytic inverse-function and stationary-branch constructions,
with positive curvature and twist, give common local contact graphs,
active endpoint rectangles and fixed derivative bounds for the densities.
Indeed the stationary derivative is bounded away from zero at the normal
orbit; compactness and the uniform complex graph bounds retain a common
smaller neighborhood. Thus one may impose any fixed
$C^{2,\beta}$ bound there, $0<\beta\le1$. The physical $C^3$ and all
other hypotheses of Proposition~\ref{prop:local-stable} are included.
The family has the connected coverage and primitive-cycle property of
Theorem~\ref{thm:stat-main}; the graph and cycle pair are not inputs.

Here is the precise loss. In each comparison allow a permutation of
representatives, one common Euclidean isometry, and changes of period
representatives along a spanning tree of the recovered incidence graph.
Allow also unmarked lattice bases whose norms and inverse norms are at
most a fixed $B$. Take the infimum of
\begin{equation}\label{eq:loss}
       \max_i\|p_{C_i}-p_{\widetilde C_i}\|_{C^2(S^1)}
                            +\|L-\widetilde L\|
\end{equation}
over these bounded presentations, applying the same orthogonal map to
the lattice and all bodies. Representatives are in a bounded root-tree
gauge. Choose $B$ large enough in terms of the prior to include every
observed primitive cycle basis: its columns are bounded finite sums of
channel placements and its determinant is bounded below by $V$.
We call this the bounded-presentation loss $\mathcal E$. It asserts
closeness of actual boundary images and an unmarked period lattice, not
independent per-edge placements. The proof establishes the explicit
comparison \eqref{eq:loss}; no metric property of this infimum is needed.
The compact physical class can equivalently be specified using these
bounded presentations modulo their finite relabellings and gauges.

\subsection{From cell probabilities to density derivatives}
Choose centered rectangles $Q_0\Subset Q$ and square grids invariant
under simultaneous reversal. A histogram records the cells of $Q$,
success outside $Q$, and failure. A single phase preparation gives one
multinomial outcome, not a separate trial for every cell.

\begin{lemma}[Cell reconstruction]\label{lem:bins}
For densities uniformly bounded in $C^{2,\beta}(Q)$, cell probabilities
on a grid of spacing $h$ determine a linear reconstruction with
\begin{equation}\label{eq:bins}
     \|\widehat f-f\|_{C^2(Q_0)}
           \le C(h^\beta+\epsilon h^{-4})
\end{equation}
when each probability has error at most $\epsilon$. The same bound
holds for the distance between two physical densities whose cell
probabilities differ by at most $\epsilon$. There are $O(h^{-2})$ cells.
\end{lemma}
\begin{proof}
The averages of $1,z,z^2$ on unit cells centered at $i=-1,0,1$ form
the invertible matrix with rows $(1,i,i^2+1/12)$. Its tensor square
recovers a coordinatewise quadratic polynomial from nine cell averages.
For a $C^{2,\beta}$ function the total-degree two Taylor polynomial
has remainder $O(h^{2+\beta})$, and it is reproduced exactly by this
operator. Its derivative errors of orders $j\le2$ are
$O(h^{2+\beta-j})$ on a fixed enlarged block. A probability error
$\epsilon$ gives average-density error $\epsilon h^{-2}$ and derivative
error $O(\epsilon h^{-2-j})$. Patch the polynomials with a smooth
partition of unity whose $j$th derivatives are $O(h^{-j})$ and whose
overlap number is bounded. Leibniz's rule gives the same bounds. The
larger rectangle supplies all boundary stencils for $Q_0$.
Apply this construction to each of two densities and subtract to obtain
the comparison assertion.
\end{proof}

We will use the smaller variance of a cell indicator. If $N$ independent
preparations are made at an active time, $p_B\le Ch^2$ and Bernstein's
inequality gives, simultaneously for all cells with probability at least
$1-\delta$, the bound
\begin{equation}\label{eq:bernstein}
 |\widehat p_B-p_B|\le C(h\sqrt{r_N}+r_N),\qquad
 r_N=\frac{\log(CJN/\delta)}{N},
\end{equation}
provided the grid has at most $CN$ cells. Dependence among the indicators
of one multinomial sample does not affect a union bound. Together with
Lemma~\ref{lem:bins}, this gives
\begin{equation}\label{eq:local-noise}
           e_N=C\{h^\beta+\sqrt{r_N}h^{-3}+r_Nh^{-4}\}.
\end{equation}
This is a variance-aware refinement of the absolute-error histogram bound,
not a lower-bound or optimality assertion.

\subsection{Analytic continuation and shape synchronization}
\begin{lemma}[A uniform continuation estimate]\label{lem:continuation}
Let $p,\widetilde p$ be periodic holomorphic functions bounded on a
common strip of positive width. If they differ by at most $e$ on a
fixed positive-length real interval, then
\[
             \|p-\widetilde p\|_{C^2(S^1)}\le Ce^\theta
\]
for some $C>0$ and $0<\theta<1$ depending on the strip, bound and
interval length. The constants are uniform under translation of the
interval on the circle.
\end{lemma}
\begin{proof}
Slit a narrower periodic strip cylinder along a closed shorter interval.
The harmonic measure $\omega$ of the slit is positive in this connected
domain, equals one on the slit, and zero on the two outer boundary
circles. Its minimum on a further narrower closed cylinder, extended
by one at the slit, is some $\theta>0$. The slit endpoints are regular
Dirichlet boundary points. The subharmonic maximum principle gives
$|p-\widetilde p|\le(2M)^{1-\omega}e^\omega$. Zeros are handled by
regularization. Cauchy estimates on fixed disks in the narrower cylinder
give the two derivative bounds on the real circle. For $e=0$ use the
identity theorem. This is the classical two-constants mechanism; its
geometric loss is consistent with the analysis in \cite{Trefethen}.
\end{proof}

\begin{lemma}[Stable incidence and geometric cycles]\label{lem:sync-stable}
Suppose two physical collections in the prior class can be paired record
by record, with an independent coherent reversal on each record, so their
active densities have $C^2$ discrepancy at most $e$. For sufficiently
small $e$ they have isomorphic incidence multigraphs and the same number
of obstacle classes. In a common root-tree gauge their corresponding
representative images, cycle displacements and calibrated covolumes
have error at most $Ce^\theta$.
\end{lemma}
\begin{proof}
Proposition~\ref{prop:local-stable} gives $Ce$ error for each pair of
actual arcs in its contact frame. Positive curvature converts them to
support functions on a common normal-angle interval. The normal-angle
map has derivative equal to curvature, bounded below; its inverse and
the support function through two derivatives depend Lipschitzly on the
arc, using the physical $C^3$ prior. Lemma~\ref{lem:continuation} gives
$\Delta=Ce^\theta$ error for each complete pair image.

Steiner centering is a bounded linear operation on support functions.
If two occurrences represent the same shape in one table, their partners
in the other have congruence distance at most $C\Delta$. They must belong
to the same class by \eqref{eq:shape-margin}. Conversely two distinct
classes cannot become congruent when $C\Delta<\nu/4$. Thus the equality
pattern of all occurrences, and hence the incidence multigraph, agrees.

Two approximate matchings of one source image to a placed representative
differ by an orthogonal transformation moving its centered support
function by $O(\Delta)$ in $L^2$. The reflection margin excludes the
incorrect component; the rotation margin bounds its angle by
$C\Delta/\eta$. Uniform third derivatives of support functions control
the induced $C^2$ error, and Steiner points control translations.
Induction along a fixed spanning tree aligns all representative images
with error $C\Delta$. The same argument for a non-tree pair and
\eqref{eq:cycle-displacement} give $C\Delta$ error in each $d_e$.
Finally the area error in the local inverse is $Ce$, and
\eqref{eq:support-area} bounds each obstacle-area error by $C\Delta$.
Summing each incidence class once proves the covolume estimate.
\end{proof}

\begin{lemma}[Integer locking for physical candidates]\label{lem:integer-lock}
Let two physical candidates have corresponding independent cycle matrices
$D,\widetilde D$ and covolumes $V,\widetilde V$ in the uniform bounded
class. If these quantities are sufficiently close, their calibrated
indices are the same integer $n$. If $n=1$, $D$ and $\widetilde D$
are bases of the two period lattices. More generally, corresponding
additional cycle vectors satisfy
\[
          nD^{-1}d_e=n\widetilde D^{-1}\widetilde d_e\in\Z^2
\]
when their discrepancy is sufficiently small, with the threshold also
depending on the bounded index.
\end{lemma}
\begin{proof}
The determinant-to-covolume ratio is Lipschitz on the bounded set with
$V$ bounded below. For physical candidates both ratios are positive
integers. A difference less than $1/2$ forces equality. Index one means
that the cycle subgroup already is the period lattice. At general index
$n$, the finite quotient $\Lambda/D\Z^2$ has order $n$, so
$n\Lambda\subset D\Z^2$. Thus the displayed vectors are integer.
Their difference is small by the inverse-matrix identity and the stated
bounds. An $\ell^\infty$ difference less than $1/2$ forces equality.
\end{proof}

The lemma applies to exact physical candidates compatible with uncertain
observations. It is not an instruction to take the integer span of noisy
real vectors, which need not be discrete, or to round uncalibrated real
numbers without an error bound. This distinction is important when no
integer copy labels are supplied.

\subsection{The pilot, estimator and preparation budget}
\begin{proof}[Proof of Theorem~\ref{thm:stat-main}]
Uniform local geometry gives $d_*,c,C>0$ such that
$cd^2\le\mu_{T_0+d}(\R^2)\le Cd^2$ for $0<d<d_*$, and zero mass
before the selected onset. One way to see the uniform quadratic order
is to scale the positive quadratic minimum of $W-T_0$ by $\sqrt d$
in \eqref{eq:density}; the two endpoint dimensions contribute $d$ and
the residual factor another $d$. Smooth Morse coordinates and the
positive twist make the two-sided bounds uniform.

Choose once and for all $\lambda>0$ with
$\sqrt{4\lambda/c}<d_*/32$, and a grid spacing less than $d_*/32$.
Query every onset bracket on this grid, including a short collar beyond
its endpoint, and let $\tau$ be the first reported mass exceeding
$2\lambda$. If the pilot errors are less than $\lambda/4$, no pre-onset
point can trigger; a grid point at offset between
$\sqrt{4\lambda/c}$ and that number plus one grid step must trigger.
Hence $0<\tau-T_0<d_*/16$. This uses no global monotonicity of the
three-impact probability. Use
$T_1=\tau+d_*/4$ and $T_2=\tau+d_*/2$. A uniform small rectangle
with $W-T_0\le d_*/8$ is strictly active at both times. All local
margins are positive. The pilot has a fixed prior-dependent number of
times per record, not a growing number as $N$ increases.

At each pilot time use $N$ fresh preparations, and at the two chosen
histogram times use fresh groups of $N$ preparations. Conditional on the
pilot, the histogram times are fixed and \eqref{eq:bernstein} applies.
Bernoulli concentration at the pilot and a union bound over all stages
and cells give pilot error $C\sqrt{r_N}$ and cell error
$C(h\sqrt{r_N}+r_N)$ with total failure probability at most $\delta$.
All preparations, including no-event outcomes, are charged.

For precision define an estimator over the compact physical prior, not
over arbitrary fitted density functions. Minimize the maximum discrepancy
from the retained pilot and histogram probabilities, divided by their
respective error tolerances, allowing a permutation of records and one
common reversal per record. Take an approximate minimizer from a fixed
countable dense family, within one of the infimum, using a first-eligible
index rule. The rule is Borel. The probability functions are continuous:
outside phase-null grazing, cell-boundary, patch-boundary and terminal-time
sets, collision records and intrinsic marks vary continuously. Uniform
separation bounds the number of collisions in a bounded queried interval,
and a fixed cell parametrization with $A$ bounded below permits dominated
convergence. Uniformly isolated contact branches also vary continuously.
Finite permutations and signs cause no obstruction. This justifies the
dense-family construction; no efficient search over analytic tables is
claimed.

On the concentration event the true table is admissible with normalized
discrepancy at most one. Thus the selected physical candidate has errors
bounded by fixed multiples of the stated tolerances. Keeping the pilot
in the discrepancy is essential: its first crossing also locates the
candidate's onset within the same fixed collar. For large $N$, candidate
pilot errors are still below $\lambda/4$, so both histograms are active
for it as well as for the truth. Lemma~\ref{lem:bins} gives the physical
pair discrepancy \eqref{eq:local-noise} with a changed constant.

Lemma~\ref{lem:sync-stable} recovers the incidence isomorphism and gives
$Ce_N^\theta$ errors in representatives, cycles and $V$. Choose, for the
comparison proof, a true spanning tree and cycle pair with calibrated
index one. The corresponding pair in the candidate has index one by
Lemma~\ref{lem:integer-lock}. Hence its cycle matrix is a period basis,
and \eqref{eq:loss} is at most $Ce_N^\theta$. The estimator was not
supplied this choice; the argument applies to every compatible candidate.

Finally take a centered square grid of spacing comparable to
$h=r_N^{1/(2\beta+6)}$. Then it has at most $CN$ cells and
\[
 h^\beta= r_N^{\beta/(2\beta+6)},\quad
 \sqrt{r_N}h^{-3}= r_N^{\beta/(2\beta+6)},\quad
 r_Nh^{-4}= r_N^{(2\beta+2)/(2\beta+6)}.
\]
The last term is smaller for $r_N<1$. There are at most $CJ$ queried
times, so the total is at most $CJN$ preparations. This proves
\eqref{eq:rate-main} and all discrete-stability assertions.
\end{proof}

The exponent is conditional on the analytic continuation geometry.
Near congruent shape classes or symmetric bodies, the incidence or
orientation thresholds deteriorate; without a uniform analytic strip,
exact image determination alone supplies no uniform noisy estimate.
At nonprimitive index the exact finite list remains valid, but a unique
estimate cannot be asserted merely by fitting arbitrarily fine histograms,
as Proposition~\ref{prop:ambiguity} demonstrates.

\section{Information, comparisons and retained results}
\label{sec:comparison}
\subsection{Boundary distance, lens data and obstacle travelling times}
The exact equivalence in Proposition~\ref{prop:lens} identifies the
information category of the local result. On a physical active rectangle,
two endpoint density functions determine a travel-time generating function
and a volume normalization; conversely these two objects determine the
densities. The affine elimination is elementary. The geometric content
lies in recovering two unknown reflecting arcs from the diagonal Hessian,
and the new global content lies in recovering incidence and the unmarked
period lattice from an unlabelled collection. We do not claim that endpoint
travel times, their tangential derivatives, or the passage to a scattering
relation are new inverse-geometric concepts.

Stefanov, Uhlmann and Vasy~\cite{SUV}, Theorem~1.1, recover a Riemannian
metric locally near a strictly convex boundary point from boundary-distance
data in dimension at least three, modulo a boundary-fixing diffeomorphism.
Their Theorem~1.3 gives global lens rigidity under a convex-foliation
hypothesis. The unknown there is a metric on a manifold with identified
boundary. Here the ambient metric is Euclidean and known, while both
reflecting curves are unknown; the measured coordinates are intrinsic
arclengths on a selected unknown source obstacle. The diagonal action
is $W(s,s)=2\ell(s)>0$, rather than the zero diagonal of an ordinary
boundary-distance function. The Schur complement in Lemma~\ref{lem:curvature}
uses a stationary interior reflection. Thus the stated metric-rigidity
theorems do not directly give our two-arc formulas; nor does our selected
planar return experiment imply their much more general metric recovery.

Gurfinkel, Noakes and Stoyanov~\cite{GNS} consider travelling times between
points on a known enclosing boundary in a curved ambient space. Their
Theorem~1 establishes conjugacy of the nontrapping scattering flows from
almost-equal travelling times. Their Theorem~2 determines finite unions of
strictly convex obstacles in a two-dimensional nonpositively curved
manifold from those travelling times. Their Lemma~7 makes explicit how
endpoint derivatives recover the incident and outgoing directions.
Noakes and Stoyanov~\cite{NS}, Theorem~1 and Sections~3--4, likewise relate
travelling-time data to conjugacy and give a constructive reconstruction
for two convex planar obstacles. These are close comparisons, not merely
background spectral references. They use families of externally observed
scattering trajectories; our data select a local source--opposite--source
return and record intrinsic endpoints on an unknown reflecting component.
The local origins and selected branch remain marks. No theorem here
converts our finite channel collection into their full exterior data.

The volume normalization has a similarly established context. Stoyanov's
extension of Santal\'o's formula~\cite{Santalo} relates phase volume and
travelling-time averages and applies to obstacle-volume recovery in the
presence of trapping. Our identity for $A$ is instead a local consequence
of the unconditional residual-time density. Its use in
$V=A+\sum_i\area(C_i)$ converts geometric cycle displacements into
an integer period index. It is this combination, not a general claim
of first recovering volume from dynamics, that enters the descent theorem.

Marked-length results have a different information set. De Simoi, Kaloshin
and Leguil~\cite{DSKL} determine analytic dispersing configurations under
axial symmetry and genericity assumptions from marked lengths. Finamore
and Leguil~\cite{FL} prove isometric rigidity of finite-horizon Sinai
billiards from an enriched marked length spectrum. We neither identify
our endpoint laws with either spectrum nor infer a stronger spectral
rigidity theorem. The present no-finite-horizon assertion concerns the
selected local channels. The generic shape-separation and primitive-cycle
hypotheses are conditions of our unlabelled reconstruction, not conclusions
about arbitrary dispersing tables. These comparisons delimit the results;
they are not an exhaustive priority classification.

\subsection{What is removed, and what is retained}
Theorem~\ref{thm:descent-main} removes the supplied obstacle identifications,
incidence graph, cross-channel contact offsets, relative edge signs and
integer lattice-copy labels on the stated shape-separated class. The
record itself still pairs its two times, specifies one local return
orientation, and uses a marked contact origin and a coherent arclength
reversal. The finite theorem also uses onset brackets and uniform priors.
``Unlabelled'' therefore describes the relations \emph{between} records,
not the absence of all experimental selection.

The area calibration is indispensable to the arithmetic reduction.
Without it, containing two independent geometric periods does not bound
the index of a finer lattice. With it, the finite list is exact, and a
primitive determinant is a directly verifiable sufficient certificate
of uniqueness. Proposition~\ref{prop:ambiguity} shows that even the raw
normalization and rank-two geometry need not give uniqueness at larger
index. This obstruction persists for exact physical density functions;
it is not a statistical artifact or an inference from formal series.

The local Lipschitz estimate is conditional on the physical bounds in
Proposition~\ref{prop:local-stable}. The global exponent is additionally
conditional on the analytic strip and continuation geometry, as in
\cite{Trefethen}, and the separation margins in Section~\ref{sec:stability}.
The integer-locking argument is a comparison of physical candidates with
known integrality, not the formation of an integer span from noisy real
vectors. None of these estimates establishes a minimax exponent.

\subsection{Comparison with the two retained regularizations}
The complete preceding article is Supplement R. Its moment-compressed
experiment uses masses, signed first moments and quadratic moments.
For a contact-jet order $K$ its whole-table bound is
\[
 C\left(A_K\epsilon^{1/(\lceil K/2\rceil+2)}+Bq^K\right)^{\vartheta},
                       \qquad 0<q,\vartheta<1.
\]
The constants $A_K$ are finite at each order but need not be uniformly
bounded. A slowly increasing regularization order gives consistency.
That theorem remains useful when the acquisition stores moment means
rather than full endpoint histograms; the exact-germ version also exposes
the even and odd triangular blocks coefficient by coefficient.

Supplement R's full-law experiment instead gives
$C(h^\beta+\epsilon h^{-4})^\vartheta$ for absolute cell-probability
accuracy $\epsilon$. It has no growing jet order, but each of its two
windows is a growing two-dimensional histogram. The present paper retains
that deterministic local estimate. Using the bound $p_B\le Ch^2$ on
cell probabilities improves the preparation accounting to
\eqref{eq:rate-main}. Simultaneously, shape separation and period-index
locking permit recovery of global combinatorial data previously supplied.
The power in the new phase-preparation bound and the power in the old
absolute-mean-error bound refer to different error parameters. They are
not competing minimax exponents on a common unqualified sensor model.

The full-law route is the primary article because it reconstructs arcs
as functions before continuation and supports the unlabelled descent.
The complete moment inverse, its dimension-sharp finite-jet designs,
registered rigidity for symmetric tables, and the old order-dependent
regularization are supplied unchanged in Supplement R. Supplement S
retains the smooth relative boundary law and its Volterra and Abel
proofs. This separation changes the reading order, not the mathematical
content or status of those results.

Finally, the retained count-data fibers are finite, formal, or smooth
Taylor-series statements in their specified categories. Neither the
full-law inverse nor the period ambiguity theorem implies an exact
analytic count-only rigidity or nonrigidity theorem. In particular,
equality of all smooth Taylor coefficients is not used to identify
smooth functions, and a formal inverse series is not assumed to converge.

\subsection*{Supplementary volumes}
Supplement R is the complete reviewed preceding article, supplied with
all of its mathematical inputs unchanged. Its moment-compressed inverse,
registered symmetric-table theorem, order-dependent regularization,
count-fiber statements and physical finite-jet experiments remain part
of this submission. Supplement S contains the complete smooth relative,
Volterra and Abel theory and its auxiliary document. Neither supplement
is cited as an independently published or accepted article. The present
primary proof is self-contained. The source index gives the precise
reading map and immutable identities of all retained material.
\subsection*{Acknowledgment and assisted analysis}
AI-assisted analysis contributed to the reconstruction, finite-index
descent and stability arguments. The named author is responsible for
checking the mathematics and references before submission. Algebraic
and numerical diagnostics and compilation are reproducibility evidence,
not formal proof certificates or independent referee acceptance.
\begin{thebibliography}{99}
\raggedright
\bibitem{DSKL}
J. De Simoi, V. Kaloshin, and M. Leguil,
\emph{Marked length spectral determination of analytic chaotic billiards
with axial symmetries}, Invent. Math. \textbf{233} (2023), 829--901.
\href{https://doi.org/10.1007/s00222-023-01191-8}{doi:10.1007/s00222-023-01191-8}.
\bibitem{FL}
D. Finamore and M. Leguil,
\emph{A CAT(0)-approach to the marked length spectral rigidity of Sinai
billiards}, arXiv:2510.18983v1 (2025).
\bibitem{GNS}
T. Gurfinkel, L. Noakes, and L. Stoyanov,
\emph{Travelling times in scattering by obstacles in curved space},
J. Differential Equations \textbf{269} (2020), 9508--9530.
\href{https://doi.org/10.1016/j.jde.2020.05.049}{doi:10.1016/j.jde.2020.05.049}.
\bibitem{NS}
L. Noakes and L. Stoyanov,
\emph{Travelling times in scattering by obstacles},
J. Math. Anal. Appl. \textbf{430} (2015), 703--717.
\href{https://doi.org/10.1016/j.jmaa.2015.05.013}{doi:10.1016/j.jmaa.2015.05.013}.
\bibitem{SUV}
P. Stefanov, G. Uhlmann, and A. Vasy,
\emph{Local and global boundary rigidity and the geodesic X-ray transform
in the normal gauge}, Ann. of Math. (2) \textbf{194} (2021), 1--95.
\href{https://doi.org/10.4007/annals.2021.194.1.1}{doi:10.4007/annals.2021.194.1.1}.
\bibitem{Santalo}
L. Stoyanov,
\emph{Santal\'o's formula and stability of trapping sets of positive measure},
J. Differential Equations \textbf{263} (2017), 2991--3008.
\href{https://doi.org/10.1016/j.jde.2017.04.019}{doi:10.1016/j.jde.2017.04.019}.
\bibitem{Trefethen}
L. N. Trefethen,
\emph{Quantifying the ill-conditioning of analytic continuation},
BIT Numer. Math. \textbf{60} (2020), 901--915.
\href{https://doi.org/10.1007/s10543-020-00802-7}{doi:10.1007/s10543-020-00802-7}.
\end{thebibliography}

\end{document}
